\begin{circuitikz}[
    american,
    voltage dir=RP,
    >=latex,
    %>={Latex[length=5mm, width=2mm]},
    alt={}
    ]

    % Define the number of rungs in the diagram
    \def\numRungs{3}

    % Calculate the length of the vertical power rails dynamically based on number of rungs
    \pgfmathsetmacro{\railLength}{-2 * \numRungs}

    %======================================================================================================
    % Component Grid
    % Spacing = 3x, 2y
    %======================================================================================================
    \foreach \i in {0,...,5} {
        \foreach \j in {0,...,20} { % Change to 0,...,19 if you need 20 individual Y points down to -40
        
        % Calculate spatial positions
        \pgfmathsetmacro{\px}{\i * 3}
        \pgfmathsetmacro{\py}{-1 - (\j * 2)}
        
        % Name coordinate using loop indices: (C-column-row) e.g., (C-0-0), (C-5-10)
        \coordinate (C-\i-\j) at (\px, \py);
        
        % OPTIONAL: Uncomment the 2 lines below to display grid labels while building
        % \fill[red] (C-\i-\j) circle (1.5pt);
        %\node[anchor=south west, scale=0.5, gray] at (C-\i-\j) {(\i,\j)};
        }
    }

    \begin{pgfonlayer}{ft}
        
        \draw[
            BakerBlack,
            very thick,
            name path=L1
        ]
        (0,0)
        -- (0,\railLength);

        \draw[
            BakerBlack,
            very thick,
            name path=L2
        ]
        (15,0)
        -- (15,\railLength);

    \end{pgfonlayer}

    \begin{pgfonlayer}{main}

        %=================================================================
        % Rung 1
        %=================================================================
        \draw[
            BakerBlack
        ]
        (C-0-0)
        to[C, color=BakerBlack, l={{\small Photo1}}] (C-1-0)
        -- (C-3-0)
        pic (CTU) {TON};

        \draw[
            BakerBlack
        ]
        (CTU-OUT)
        -- (C-5-0);

        \node[text=BakerBlack, anchor=south] at (CTU-TOP) {{\small CTU}};
        \node[text=BakerBlack, anchor=north] at (CTU-TOP) {{\small PartCount}};
        \node[text=BakerBlack, anchor=west] at (CTU-PRE) {{\small 10}};
        \node[text=BakerBlack, anchor=west] at (CTU-ACC) {{\small 9}};

        %=================================================================
        % Rung 3 (timer takes 2 rungs)
        %=================================================================
        \draw[
            BakerBlack
        ]
        (C-0-2)
        to[C, color=BakerBlack, l={{\small BatchDone}}] (C-1-2)
        -- (C-4-2)
        node[draw=BakerBlack, text=BakerBlack, font=\small, minimum width=1cm, minimum height=0.5cm, anchor=west, rounded corners=0.25cm](res){RES};

        \node[
            above,
            text=BakerBlack
        ]
        at (res.north)
        {{\small PartCount}};

        \draw[
            BakerBlack
        ]
        (res.east)
        -- (C-5-2);

    \end{pgfonlayer}

    %=================================================================
    % Background layer - highlight any logic that should be true
    % UpnorthGreen, opacity=0.4, line width=10pt
    %=================================================================
    \begin{pgfonlayer}{bg}
        
        %=================================================================
        % Rung 1
        %=================================================================
        % \draw[
        %     UpNorthGreen,
        %     opacity=0.4,
        %     line width=10pt,
        % ]
        % (C-0-0)
        % -- (C-1-0);

        % \draw[
        %     UpNorthGreen,
        %     opacity=0.4,
        %     line width=10pt,
        % ]
        % (CTU-EN1)
        % -- (CTU-EN2);

    \end{pgfonlayer}
    
\end{circuitikz}